\documentclass[12pt]{article}
\usepackage[utf8]{inputenc}
\usepackage[T1]{fontenc}
\usepackage{hyperref}
\title{Peruvian Sign Language Language Resources for K-12 Curriculum Development}
\author{Compiled by Claude Sonnet 5.5 (claude-sonnet-5-5) for Andr\'as Kornai}
\date{}
\begin{document}
\maketitle
\section{Introduction}
Peruvian Sign Language (Lengua de Se\~nas Peruana) is recognised by Law 29535. Searches found Ministry of Education publications and legal documents; no official dictionary was found.
\section{Primary Educational Resources}
\begin{itemize}
\item Manual de lenguaje de se\~nas peruanas / Per\'u Mama Makillawan Rimaykuna (Minedu): \url{https://repositorio.minedu.gob.pe/handle/20.500.12799/5544}
\item Lengua de se\~nas peruana: gu\'ia para el aprendizaje, vocabulario b\'asico (Minedu): \url{https://repositorio.minedu.gob.pe/handle/20.500.12799/5545}
\end{itemize}
\section{Linguistic and Reference Materials}
No academic grammar was found in these searches.
\section{Religious and Community Resources}
None found.
\section{Digital and Audio Resources}
No video collections were identified in the results.
\section{Curriculum Development Context}
Legal and policy background (Congress of Peru briefing on sign language): \url{https://www.congreso.gob.pe/Docs/DGP/DIDP/files/nir_2024-2025/nir-35-lengua-de-se\~nas.pdf}. News on the regulation: \url{https://andina.pe/ingles/noticia-nuevo-reglamento-sobre-lengua-senas-beneficia-a-mas-700000-personas-679001.aspx}. The regulation requires Minedu to guarantee deaf students' education and to approve profiles for interpreters and linguistic models.
\end{document}